\documentclass[10pt,a4paper]{amsart}
\usepackage{amsmath,amssymb,amsthm,graphicx}
\usepackage{verbatim}
\newenvironment{reference}{\comment}{\endcomment}
%\newenvironment{reference}{}{}
\newenvironment{slogan}{\comment}{\endcomment}
\newenvironment{history}{\comment}{\endcomment}

% For commutative diagrams we use Xy-pic
\usepackage[all]{xy}

% We use 2cell for 2-commutative diagrams.
\xyoption{2cell}
% Theorem environments.
%
\theoremstyle{plain}
\newtheorem{theorem}[subsection]{Theorem}
\newtheorem{proposition}[subsection]{Proposition}
\newtheorem{lemma}[subsection]{Lemma}

\theoremstyle{definition}
\newtheorem{definition}[subsection]{Definition}
\newtheorem{example}[subsection]{Example}
\newtheorem{exercise}[subsection]{Exercise}
\newtheorem{situation}[subsection]{Situation}

\theoremstyle{remark}
\newtheorem{remark}[subsection]{Remark}
\newtheorem{remarks}[subsection]{Remarks}

\numberwithin{equation}{subsection}

% Macros
%
\def\lim{\mathop{\mathrm{lim}}\nolimits}
\def\colim{\mathop{\mathrm{colim}}\nolimits}
\def\Spec{\mathop{\mathrm{Spec}}}
\def\Hom{\mathop{\mathrm{Hom}}\nolimits}
\def\Ext{\mathop{\mathrm{Ext}}\nolimits}
\def\SheafHom{\mathop{\mathcal{H}\!\mathit{om}}\nolimits}
\def\SheafExt{\mathop{\mathcal{E}\!\mathit{xt}}\nolimits}
\def\Sch{\mathit{Sch}}
\def\Mor{\mathop{\mathrm{Mor}}\nolimits}
\def\Ob{\mathop{\mathrm{Ob}}\nolimits}
\def\Sh{\mathop{\mathit{Sh}}\nolimits}
\def\NL{\mathop{N\!L}\nolimits}
\def\CH{\mathop{\mathrm{CH}}\nolimits}
\def\proetale{{pro\text{-}\acute{e}tale}}
\def\etale{{\acute{e}tale}}
\def\QCoh{\mathit{QCoh}}
\def\Ker{\mathop{\mathrm{Ker}}}
\def\Im{\mathop{\mathrm{Im}}}
\def\Coker{\mathop{\mathrm{Coker}}}
\def\Coim{\mathop{\mathrm{Coim}}}

% Boxtimes
%
\DeclareMathSymbol{\boxtimes}{\mathbin}{AMSa}{"02}

%
% Macros for moduli stacks/spaces
%
\def\QCohstack{\mathcal{QC}\!\mathit{oh}}
\def\Cohstack{\mathcal{C}\!\mathit{oh}}
\def\Spacesstack{\mathcal{S}\!\mathit{paces}}
\def\Quotfunctor{\mathrm{Quot}}
\def\Hilbfunctor{\mathrm{Hilb}}
\def\Curvesstack{\mathcal{C}\!\mathit{urves}}
\def\Polarizedstack{\mathcal{P}\!\mathit{olarized}}
\def\Complexesstack{\mathcal{C}\!\mathit{omplexes}}
% \Pic is the operator that assigns to X its picard group, usage \Pic(X)
% \Picardstack_{X/B} denotes the Picard stack of X over B
% \Picardfunctor_{X/B} denotes the Picard functor of X over B
\def\Pic{\mathop{\mathrm{Pic}}\nolimits}
\def\Picardstack{\mathcal{P}\!\mathit{ic}}
\def\Picardfunctor{\mathrm{Pic}}
\def\Deformationcategory{\mathcal{D}\!\mathit{ef}}
\usepackage{cleveref}
\crefname{originallemma}{original lemma}{original lemmas}
\Crefname{originallemma}{Original lemma}{Original lemmas}
\makeatletter
\newlabel{lemma-compare-cohomology-point}{{03QU}{source}{}{sourcepointer-1}{}}
\newlabel{lemma-compare-cohomology-point@cref}{{[originallemma][0][]03QU}{[1][0][]source}{}{}{}}
\makeatother
\crefname{originalproposition}{original proposition}{original propositions}
\Crefname{originalproposition}{Original proposition}{Original propositions}
\makeatletter
\newlabel{proposition-serre-galois}{{03R8}{source}{}{sourcepointer-2}{}}
\newlabel{proposition-serre-galois@cref}{{[originalproposition][0][]03R8}{[1][0][]source}{}{}{}}
\makeatother
\crefname{originallemma}{original lemma}{original lemmas}
\Crefname{originallemma}{Original lemma}{Original lemmas}
\makeatletter
\newlabel{lemma-curve-brauer-zero}{{03RF}{source}{}{sourcepointer-3}{}}
\newlabel{lemma-curve-brauer-zero@cref}{{[originallemma][0][]03RF}{[1][0][]source}{}{}{}}
\makeatother
\title{Vanishing of the etale cohomology of the multiplicative group over a function field of a curve}
\author{The Stacks Project Authors}
\date{stacks-project:etale-cohomology.tex:local-snapshot-20261011; GNU FDL 1.2 or later}
\begin{document}
\maketitle
\noindent\textit{Editorial note.} All delivered mathematical text is the author's own, copied from the pinned Stacks Project chapter file etale-cohomology.tex; no statement or proof step was rewritten, reordered or removed.\par\medskip
\noindent\textit{Editorial note.} The proof is the author's complete argument; both reductions are inside the counted span.\par\medskip
\noindent\textit{Editorial note.} Every labelled reference occurring inside the counted spans is carried in assembly.external\_references as a non-load-bearing navigation pointer (same-chapter ones at their definition site, cross-chapter ones at the invocation inside the counted span), so no load-bearing mathematics is replaced by a stub. No bibliography citation occurs inside the counted spans.\par\medskip
\noindent\textit{Editorial note.} The checked chapter list for same-name collisions is the set of stacks chapters staged in this repository: chow, cohomology, divisors, etale-cohomology, flat, more-algebra, more-morphisms, morphisms, spaces-more-morphisms, varieties together with this card's chapter etale-cohomology.\par\medskip
\noindent\textit{Editorial note.} The preamble is the chapter's own native preamble excerpt already used for this source: the Xy-pic setup (native preamble lines 14-18) and the native theorem declarations with the mathematical macros (native preamble lines 154-217).\par\medskip
\noindent\textit{Allowed standard prerequisites (editorial).}\begin{itemize}
\item Comparison of Galois cohomology with etale cohomology of the spectrum of a field. Same chapter, the identification of H\textasciicircum{}q\_etale(Spec(K), G\_m) with the Galois cohomology of the multiplicative group of the separable closure. Supplies the reduction to Galois cohomology.
\item Serre's vanishing theorem under the vanishing of Brauer groups of finite extensions. Same chapter, the proposition deducing the vanishing of the Galois cohomology from the vanishing of the Brauer groups of all finite extensions, together with the vanishing of the Brauer group of a function field of a curve. Supplies the two ingredients of the proof.
\end{itemize}\medskip
\ifdefined\hypertarget\hypertarget{sourcepointer-1}{}\fi
\noindent\textit{Source reference.} Original lemma 03QU: sources/stacks-etale-cohomology/pinned-original-source.tex, source lines 9100--9100. Navigation pointer invoked at proof line 10916; the counted span contains the pointer only and the cited result is not reproduced inside it, so the pointer is registered as a non-load-bearing navigation stub in the established form for Stacks material. No bibliography citation is involved.\par\medskip
\ifdefined\hypertarget\hypertarget{sourcepointer-2}{}\fi
\noindent\textit{Source reference.} Original proposition 03R8: sources/stacks-etale-cohomology/pinned-original-source.tex, source lines 10694--10694. Navigation pointer invoked at proof line 10917; the counted span contains the pointer only and the cited result is not reproduced inside it, so the pointer is registered as a non-load-bearing navigation stub in the established form for Stacks material. No bibliography citation is involved.\par\medskip
\ifdefined\hypertarget\hypertarget{sourcepointer-3}{}\fi
\noindent\textit{Source reference.} Original lemma 03RF: sources/stacks-etale-cohomology/pinned-original-source.tex, source lines 10894--10894. Navigation pointer invoked at proof line 10926; the counted span contains the pointer only and the cited result is not reproduced inside it, so the pointer is registered as a non-load-bearing navigation stub in the established form for Stacks material. No bibliography citation is involved.\par\medskip
\begin{lemma}
\label{lemma-cohomology-Gm-function-field-curve}
Let $k$ be an algebraically closed field and $K/k$ a field extension
of transcendence degree 1. Then for all $q \geq 1$,
$H_\etale^q(\Spec(K), \mathbf{G}_m) = 0$.
\end{lemma}

\begin{proof}
Recall that
$H_\etale^q(\Spec(K), \mathbf{G}_m) = H^q(\text{Gal}(K^{sep}/K), (K^{sep})^*)$
by Lemma \ref{lemma-compare-cohomology-point}.
Thus by Proposition \ref{proposition-serre-galois}
it suffices to show that if $K'/K$ is a finite field extension, then
$\text{Br}(K') = 0$. Now observe that $K' = \colim K''$, where $K''$ runs
over the finitely generated subextensions of $k$ contained in $K'$ of
transcendence degree $1$.
Note that $\text{Br}(K') = \colim \text{Br}(K'')$ which reduces us
to a finitely generated field extension $K''/k$ of transcendence
degree $1$. Such a field is the function field of a curve over $k$,
hence has trivial Brauer group by
Lemma \ref{lemma-curve-brauer-zero}.
\end{proof}

\end{document}
